\chapter[Phân loại phản hồi của sinh viên]{Phân loại phản hồi của sinh viên theo cảm xúc và chủ đề}

\section{Bài toán và hướng tiếp cận}

Mỗi học kỳ, Trường Đại học Công nghệ Thông tin (ĐHQG-HCM) thu thập ý kiến phản hồi tự do của sinh viên về giảng viên,
chương trình đào tạo và cơ sở vật chất. Với hơn 16 000 câu phản hồi, việc đọc và phân loại thủ công tốn nhiều thời
gian. Bài tập này xây dựng mô hình tự động gán cho mỗi câu phản hồi hai nhãn:

\begin{itemize}
  \item \textbf{Cảm xúc} (\texttt{sentiment}, nhiệm vụ chính): \texttt{negative}, \texttt{neutral}, \texttt{positive}.
  \item \textbf{Chủ đề} (\texttt{topic}, nhiệm vụ phụ): \texttt{lecturer} (giảng viên), \texttt{program} (chương trình
        đào tạo), \texttt{facility} (cơ sở vật chất), \texttt{others} (chủ đề khác).
\end{itemize}

Đây là bài toán phân tích cảm xúc, một dạng phân loại văn bản có giám sát. Vì mô hình phân loại chỉ làm việc với số,
mỗi câu được biểu diễn thành một vector. Quy trình gồm các bước: kiểm tra chất lượng và phân tích khám phá dữ liệu; tiền
xử lý văn bản (làm sạch, đổi viết tắt, tách từ tiếng Việt); so sánh ba bộ phân lớp của scikit-learn trên đặc trưng
TF-IDF; thay TF-IDF bằng embedding câu tiếng Việt cho bộ phân lớp tốt nhất; đánh giá một lần trên tập test và phân tích
lỗi. Notebook chạy được trên Google Colab với \textit{Runtime > Run all}.

\section{Bộ dữ liệu}

\begin{itemize}
  \item \textbf{Nguồn:} bộ dữ liệu UIT-VSFC (Vietnamese Students' Feedback Corpus) \cite{text-vsfc}, tải từ trang Hugging Face
        của nhóm tác giả (\texttt{uitnlp/vietnamese\_students\_feedback}), không cần đăng nhập.
  \item \textbf{Giấy phép:} dùng cho mục đích nghiên cứu theo thỏa thuận của nhóm tác giả. Nhóm trích dẫn nguồn và không
        đưa dữ liệu lên repo.
  \item \textbf{Nội dung:} câu phản hồi trong khảo sát cuối học kỳ, từ năm học 2014-2015 đến 2016-2017, được gán nhãn thủ
        công. Văn bản đã được viết thường, tách dấu câu, ẩn danh tên giảng viên (\texttt{wzjwz<số>}) và chuẩn hóa viết tắt
        trước khi công bố.
  \item \textbf{Quy tắc gán nhãn:} \texttt{negative} gồm cả lời phàn nàn lẫn đề nghị cải thiện; \texttt{neutral} là câu
        không nêu ý kiến hoặc chưa trọn ý, kể cả lời cảm ơn; câu vừa khen vừa chê được gán theo ý mạnh hơn.
  \item \textbf{Độ tin cậy của nhãn:} mức đồng thuận giữa người gán nhãn là 91.2\% với cảm xúc nhưng chỉ 71.07\% với chủ
        đề, nên nhãn chủ đề kém chắc chắn hơn.
\end{itemize}

Nhóm giữ nguyên cách chia train / dev / test của tác giả và dùng làm train / val / test, để kết quả so sánh được
với bài báo gốc (Bảng~\ref{tab:text-split}). Bộ dữ liệu đáp ứng yêu cầu của đề với dữ liệu văn bản: 16 175 câu, nhiều
hơn mức tối thiểu 2 000, và 99.9\% số câu có chữ tiếng Việt có dấu.

\begin{table}[H]\centering
\caption{Kích thước ba tập và tỉ lệ lớp \texttt{neutral}.}\label{tab:text-split}
\begin{tabular}{lrrr}
\toprule
Tập & Số câu & Tỉ lệ & Tỉ lệ \texttt{neutral} \\
\midrule
train & 11 426 & 70.6\% & 4.0\% \\
val & 1 583 & 9.8\% & 4.6\% \\
test & 3 166 & 19.6\% & 5.3\% \\
\bottomrule
\end{tabular}
\end{table}

\section{Kiểm tra chất lượng dữ liệu}

Mục này chỉ chẩn đoán vấn đề; cách xử lý được áp dụng ở bước tiền xử lý. Bảng~\ref{tab:text-quality} tổng hợp kết quả.

\begin{table}[H]\centering
\caption{Các vấn đề chất lượng dữ liệu và cách xử lý.}\label{tab:text-quality}
\begin{tabular}{>{\raggedright\arraybackslash}p{4.2cm}>{\raggedright\arraybackslash}p{4.6cm}>{\raggedright\arraybackslash}p{5.4cm}}
\toprule
Vấn đề & Kết quả kiểm tra & Cách xử lý \\
\midrule
Giá trị thiếu, câu rỗng & 0 câu & nếu có thì bỏ khỏi train/val/test \\
Unicode không ở dạng NFC & 0 câu & vẫn chuẩn hóa NFC cho câu mới \\
Câu trùng trong train & 1 câu, hai bản mang nhãn trái ngược & bỏ bản lặp khỏi train \\
Câu train có trong val/test & 0 câu & nếu có thì bỏ khỏi train \\
Câu rất ngắn / rất dài & 141 câu không quá 2 từ, 663 câu trên 30 từ & giữ lại, vì đều là câu trả lời thật \\
Tên ẩn danh \texttt{wzjwz<số>} & 211 câu, 233 mã khác nhau & thay bằng một token chung \texttt{personname} \\
Biểu tượng cảm xúc (\texttt{colonsmile}, ...) & 85 câu & giữ lại vì mang cảm xúc \\
Viết tắt / teencode & 25 token trong 21 câu & đổi sang dạng đầy đủ theo từ điển viết tắt \\
Nhãn mất cân bằng & \texttt{neutral} chỉ 4.0\% tập train & chọn mô hình theo macro-F1, dùng \texttt{class\_weight="balanced"} \\
\bottomrule
\end{tabular}
\end{table}

Từ điển viết tắt (gồm các cặp từ viết tắt và dạng đầy đủ) được xây dựng bằng cách tách
toàn bộ dữ liệu thành từ, dùng LLM rà soát để tìm từ viết tắt và cách viết sai, rồi kiểm tra thủ công từng từ trong câu
gốc. Các từ dễ hiểu sai bị loại, ví dụ \texttt{k} (phần còn lại của \texttt{k6}, \texttt{320k} sau khi bỏ chữ số) và
\texttt{bt} (có thể là ``bài tập'' hoặc ``bình thường''). Từ điển cũng có thêm các cách viết tắt phổ biến như
\texttt{gv}, \texttt{sv} để xử lý câu mới do người dùng nhập. Kết quả cho thấy dữ liệu gần như không có viết tắt, đúng
như bài báo mô tả.

\section{Phân tích khám phá dữ liệu}

Phân tích khám phá chỉ dùng tập train, trừ Hình~\ref{fig:text-labels} so sánh ba tập.

\reportfig{fig01_label_distribution.png}{Cả hai nhãn đều mất cân bằng: trong train, \texttt{neutral} chỉ chiếm 4.0\%
và chủ đề \texttt{lecturer} chiếm 71.5\%. Tỉ lệ giữa ba tập chênh nhau tối đa 2.1 điểm phần trăm.}{fig:text-labels}

Vì lớp \texttt{neutral} rất nhỏ, accuracy sẽ che mất việc mô hình nhận sai lớp này. Do đó mô hình được chọn theo
macro-F1, thước đo lấy trung bình F1 của các lớp với trọng số như nhau.

\reportfig{fig02_topic_sentiment_heatmap.png}{Tỉ lệ và số câu tiêu cực theo chủ đề: \texttt{facility} 95.4\% (474 câu),
\texttt{program} 78.0\% (1 716 câu), \texttt{others} 40.2\% (226 câu), \texttt{lecturer} 35.6\% (2 909 câu).
Phản hồi về cơ sở vật chất ít nhưng gần như toàn là phàn nàn, còn giảng viên chiếm 54.6\% tổng số câu tiêu
cực.}{fig:text-topic}

Tỉ lệ từng cặp chủ đề và cảm xúc trên tập train khớp với bảng thống kê của bài báo gốc (chênh dưới 1 điểm phần trăm ở mọi
ô), trừ ô \texttt{lecturer} x \texttt{positive} (44.38\% so với 33.57\%). Ô này có vẻ là lỗi đánh máy của bài báo: hàng
\texttt{lecturer} của bài báo chỉ cộng lại được 60.76\% thay vì 71.76\%.

\reportfig{fig03_sentence_length.png}{Trung vị độ dài là 12 từ với câu \texttt{negative}, 9 với \texttt{positive} và 6
với \texttt{neutral}. Câu tiêu cực dài hơn vì thường kèm lý do và đề xuất, nhưng các phân phối chồng lấn nhiều nên độ
dài không được dùng làm đặc trưng.}{fig:text-length}

Chia theo khoảng 5 từ như bài báo, 74.6\% số câu train không quá 15 từ, thấp hơn mức ``hơn 83\%'' của bài báo, có thể vì
bài báo đếm từ sau khi tách từ còn nhóm đếm âm tiết. Khoảng 43\% câu \texttt{others} chỉ có 1-5 từ.

\reportfig{fig04_top_ngrams.png}{Câu tích cực chủ yếu chứa từ khen (\textit{nhiệt tình}, \textit{dễ}), câu tiêu cực chứa
từ góp ý, phàn nàn (\textit{nên}, \textit{cần}, \textit{hơn}), câu trung tính chứa cụm kiểu ``không có ý kiến''. Từ vựng
tách biệt rõ giữa các lớp, nên biểu diễn câu bằng tần suất từ là hướng hợp lý.}{fig:text-ngrams}

\section{Tiền xử lý dữ liệu}

Các bước được áp dụng theo thứ tự sau. Các quy tắc theo từng dòng được áp dụng như nhau cho mọi tập, còn mọi thống kê
học từ dữ liệu (từ vựng TF-IDF, bộ mã hóa nhãn) chỉ được fit trên train; validation và test không bị bỏ dòng nào.

\begin{enumerate}
  \item Chuẩn hóa Unicode NFC và bỏ câu trùng khỏi train: 16 175 câu còn 16 174.
  \item Hàm \texttt{clean\_text()}: viết thường, thay tên ẩn danh bằng \texttt{personname}, bỏ dấu câu và chữ số, đổi
        viết tắt theo từ điển viết tắt, gộp khoảng trắng.
  \item Tách từ tiếng Việt bằng \texttt{pyvi}: nối các âm tiết của một từ bằng dấu \texttt{\_} (ví dụ \textit{giảng
        viên} thành \texttt{giảng\_viên}); tìm được 49 749 từ nhiều âm tiết. Văn bản đã tách từ dùng cho TF-IDF, văn bản
        chưa tách từ dùng cho mô hình embedding.
  \item Mã hóa nhãn bằng \texttt{LabelEncoder}.
  \item Biểu diễn TF-IDF: unigram và bigram, bỏ n-gram chỉ gặp một lần, chuẩn hóa L2 để câu dài và câu ngắn có cùng thang
        đo. Từ vựng có 13 025 đặc trưng.
\end{enumerate}

\section{Xây dựng mô hình}

Mô hình được huấn luyện trên train và chọn theo macro-F1 trên validation; tập test chỉ dùng một lần. Ba bộ phân lớp của
scikit-learn được so sánh với tham số mặc định: Naive Bayes (\texttt{MultinomialNB}), Logistic Regression và LinearSVC.
Bài báo gốc cũng dùng Logistic Regression làm baseline. Vì lớp
\texttt{neutral} rất nhỏ, Logistic Regression và LinearSVC luôn dùng \texttt{class\_weight="balanced"} để lỗi ở lớp
nhỏ bị phạt nặng hơn.

\subsection{So sánh các bộ phân lớp trên TF-IDF}

\reportfig{fig05_model_comparison.png}{Macro-F1 validation của ba bộ phân lớp trên cùng đặc trưng TF-IDF. LinearSVC tốt
nhất ở cả hai nhãn (0.773 với \texttt{sentiment}, 0.748 với \texttt{topic}).}{fig:text-models}

LinearSVC đạt macro-F1 validation cao nhất ở cả hai nhãn, Logistic Regression kém khoảng 2 điểm, còn Naive Bayes thấp
hơn hẳn (0.605 và 0.480) dù accuracy vẫn cao (0.886 với \texttt{sentiment}): mô hình này làm tốt ở lớp lớn nhưng kém ở các
lớp nhỏ như \texttt{neutral} và \texttt{others}. Đây là ví dụ cho thấy vì sao phải chọn mô hình theo macro-F1.

\subsection{Embedding tiếng Việt có giúp bộ phân lớp tốt nhất không?}

Giữ LinearSVC, nhóm thay TF-IDF bằng embedding câu của \url{AITeamVN/Vietnamese_Embedding} \cite{text-emb}, mô hình embedding câu
tiếng Việt được fine-tune từ BGE-M3 (vector 1024 chiều, giấy phép Apache 2.0). Mô hình được dùng như một hàm trích đặc
trưng có sẵn, không fine-tune thêm.

\begin{table}[H]\centering
\caption{LinearSVC trên TF-IDF và trên embedding tiếng Việt (macro-F1 validation).}\label{tab:text-emb-val}
\begin{tabular}{lrr}
\toprule
Nhãn & TF-IDF & Embedding \\
\midrule
\texttt{sentiment} & 0.7729 & \textbf{0.8244} \\
\texttt{topic} & 0.7484 & \textbf{0.7496} \\
\bottomrule
\end{tabular}
\end{table}

Với \texttt{sentiment}, embedding tốt hơn rõ rệt (+5.2 điểm); với \texttt{topic}, hai cách gần như bằng nhau. Có thể vì
cảm xúc phụ thuộc vào nghĩa của cả câu, còn chủ đề thường lộ ra qua vài từ khóa mà TF-IDF đã bắt được. Theo cách chọn
mô hình, \textbf{embedding tiếng Việt + LinearSVC} là mô hình cuối cùng cho cả hai nhãn.

\section{Đánh giá trên tập test}

\begin{table}[H]\centering
\caption{Kết quả trên tập test. Hàng in đậm là mô hình cuối cùng, chọn theo macro-F1 validation.}\label{tab:text-test}
\begin{tabular}{llrrr}
\toprule
Nhãn & Mô hình & Accuracy & Macro-F1 & Weighted-F1 \\
\midrule
\texttt{sentiment} & Naive Bayes, TF-IDF & 0.8699 & 0.5955 & 0.8466 \\
 & Logistic Regression, TF-IDF & 0.8692 & 0.7201 & 0.8748 \\
 & LinearSVC, TF-IDF & 0.8961 & 0.7388 & 0.8917 \\
 & \textbf{LinearSVC, embedding} & \textbf{0.9169} & \textbf{0.8011} & \textbf{0.9172} \\
\midrule
\texttt{topic} & Naive Bayes, TF-IDF & 0.8184 & 0.5440 & 0.7809 \\
 & Logistic Regression, TF-IDF & 0.8143 & 0.7283 & 0.8282 \\
 & LinearSVC, TF-IDF & 0.8547 & 0.7516 & 0.8562 \\
 & \textbf{LinearSVC, embedding} & \textbf{0.8500} & \textbf{0.7522} & \textbf{0.8574} \\
\bottomrule
\end{tabular}
\end{table}

Thứ tự các mô hình trên test giống trên validation ở cả hai nhãn. Weighted-F1 của mô hình cuối cùng (0.917 và 0.857) cao
hơn mô hình Logistic Regression của bài báo gốc (khoảng 0.88 và 0.84). Với \texttt{topic}, embedding và TF-IDF gần như bằng
nhau.

\reportfig{fig06_confusion_matrices.png}{Ma trận nhầm lẫn của mô hình cuối cùng (chuẩn hóa theo hàng). Recall đạt 95.0\%
với \texttt{negative} và 92.8\% với \texttt{positive} nhưng chỉ 53.3\% với \texttt{neutral}; với chủ đề, \texttt{others}
chỉ đạt 56.0\%.}{fig:text-cm}

\begin{table}[H]\centering
\caption{F1 từng lớp so với mô hình Logistic Regression của bài báo gốc trên cùng tập test.}\label{tab:text-perclass}
\begin{tabular}{llrrr}
\toprule
Nhãn & Lớp & Nhóm & Bài báo & Chênh lệch \\
\midrule
\texttt{sentiment} & \texttt{negative} & 0.9409 & 0.9052 & +0.036 \\
 & \texttt{neutral} & 0.5251 & 0.3399 & +0.185 \\
 & \texttt{positive} & 0.9374 & 0.9132 & +0.024 \\
\midrule
\texttt{topic} & \texttt{facility} & 0.9060 & 0.8873 & +0.019 \\
 & \texttt{lecturer} & 0.9149 & 0.9112 & +0.004 \\
 & \texttt{others} & 0.4635 & 0.3810 & +0.083 \\
 & \texttt{program} & 0.7245 & 0.6719 & +0.053 \\
\bottomrule
\end{tabular}
\end{table}

Mô hình cuối cùng có F1 cao hơn bài báo ở cả 7 lớp, tăng nhiều nhất ở các lớp nhỏ và khó (\texttt{neutral},
\texttt{others}, \texttt{program}). Tuy vậy \texttt{neutral} (0.53) và \texttt{others} (0.46) vẫn là hai lớp yếu nhất.

\subsection{TF-IDF hay embedding tiếng Việt?}

\reportfig{fig07_tfidf_vs_embedding.png}{Cùng bộ phân lớp LinearSVC, embedding tiếng Việt tăng macro-F1 test từ 0.7388
lên 0.8011 với \texttt{sentiment} (+6.2 điểm), còn với \texttt{topic} gần như không đổi (0.7516 và
0.7522).}{fig:text-emb}

\subsection{Mô hình dựa vào những từ nào?}

Mô hình cuối cùng dùng embedding 1024 chiều, các chiều này không ứng với từ nào nên không đọc được mô hình dựa vào từ gì.
Vì vậy nhóm phân tích mô hình TF-IDF tốt nhất (LinearSVC trên TF-IDF), mô hình có một trọng số cho mỗi cặp (lớp,
n-gram). Để kiểm chứng các n-gram mạnh nhất có thật sự quyết định dự đoán, nhóm xóa chúng khỏi vector test và so với
việc xóa cùng số lượng n-gram ngẫu nhiên.

\reportfig{fig08_top_ngrams.png}{N-gram mạnh nhất theo lớp cảm xúc của LinearSVC trên TF-IDF. \texttt{positive} dựa vào
\textit{nhiệt\_tình}, \textit{tốt}; \texttt{negative} dựa vào \textit{chưa}, \textit{cần}, \textit{nên}. Xóa 60 n-gram
mạnh nhất làm macro-F1 giảm từ 0.7388 xuống 0.6544, trong khi xóa 60 n-gram ngẫu nhiên chỉ còn 0.7373.}{fig:text-words}

Các động từ góp ý \textit{cần}, \textit{nên} là tín hiệu \texttt{negative} mạnh vì theo quy tắc gán nhãn, đề nghị cải
thiện được tính là tiêu cực. Các n-gram mạnh nhất của \texttt{neutral} (\textit{rất ít}, \textit{nhưng hơi}) mang sắc
thái chê nhẹ hoặc lưng chừng, cho thấy lớp này không có từ vựng riêng rõ ràng. Với chủ đề cơ sở vật chất, các từ đẩy
câu về tiêu cực mạnh nhất là \textit{máy\_chiếu}, \textit{phòng}, \textit{trang thiết\_bị}, \textit{wifi}.

\subsection{Mô hình sai ở đâu?}

Mô hình cảm xúc cuối cùng đoán sai 263 trên 3 166 câu test (8.3\%). Kiểu sai nhiều nhất là câu \texttt{positive} bị
đoán thành \texttt{negative} (25.5\% số lỗi), và 61.2\% số lỗi có dính đến lớp \texttt{neutral} dù lớp này chỉ có 167
câu.

\reportfig{fig09_error_analysis.png}{Tỉ lệ đoán sai theo độ dài câu, theo loại từ trong câu và theo chủ đề. Câu có từ đối
lập (\textit{nhưng}, \textit{tuy}, ...) sai 18.4\%, gấp 2.4 lần câu không có (7.8\%); câu 1-5 từ sai nhiều nhất (10.6\%);
câu thuộc chủ đề \texttt{others} sai 30.2\%.}{fig:text-errors}

Từ phủ định gần như không ảnh hưởng đến tỉ lệ sai (7.9\% so với 8.4\%), còn câu có từ góp ý lại ít sai hơn (6.1\% so với
8.8\%) vì dữ liệu luôn gán lời góp ý là \texttt{negative}. Nguồn lỗi chính là câu vừa khen vừa chê, một giới hạn của việc
gán một nhãn cho cả câu chứ không chỉ của mô hình.

\section{Dự đoán phản hồi mới}

Hàm \texttt{predict()} chạy toàn bộ quy trình trên văn bản thô: làm sạch bằng cùng hàm \texttt{clean\_text()}, mã hóa câu
bằng mô hình embedding, rồi dự đoán bằng hai mô hình cuối cùng. Câu có viết tắt như \textit{``gv dạy ko hiểu gì cả, nói
quá nhanh''} được xếp \texttt{negative}, \texttt{lecturer}; câu đề nghị \textit{``nên cho thêm bài tập thực hành cuối mỗi
chương''} được xếp \texttt{negative}, \texttt{program}, đúng quy tắc gán nhãn. Câu \textit{``thầy dạy dễ hiểu nhưng phòng
học nóng quá''} chỉ nhận một cặp nhãn (\texttt{negative}, \texttt{lecturer}); khi tách câu tại từ \textit{nhưng}, hai
mệnh đề được xếp đúng là (\texttt{positive}, \texttt{lecturer}) và (\texttt{negative}, \texttt{facility}).

\section{Kết luận}

Mô hình cuối cùng, embedding tiếng Việt kết hợp LinearSVC, đạt macro-F1 0.8011 với cảm xúc và 0.7522 với chủ đề trên tập
test, cao hơn baseline Naive Bayes (0.5955 và 0.5440) và cao hơn mô hình Logistic Regression của bài báo gốc ở mọi lớp.
Embedding tiếng Việt giúp ích rõ với nhãn cảm xúc (+6.2 điểm so với TF-IDF) nhưng gần như không thay đổi kết quả với nhãn
chủ đề.

\subsection*{Hạn chế}
\begin{itemize}
  \item Mỗi câu chỉ có một nhãn, nên câu vừa khen vừa chê không được biểu diễn đúng và là nguồn lỗi chính.
  \item Lớp \texttt{neutral} và \texttt{others} vẫn yếu (F1 khoảng 0.53 và 0.46): hai lớp nhỏ, không có từ vựng riêng và
        có một số nhãn mơ hồ.
  \item Theo quy tắc gán nhãn, lời đề nghị được tính là \texttt{negative}, nên lớp này lẫn cả phàn nàn và góp ý.
  \item Các bộ phân lớp dùng tham số mặc định, chưa được tinh chỉnh; embedding chỉ được dùng làm hàm trích đặc trưng.
  \item Dữ liệu đến từ một trường đại học và đã được chuẩn hóa, nên kết quả với phản hồi kiểu mạng xã hội có thể thấp hơn.
\end{itemize}

\subsection*{Hướng mở rộng}
Fine-tune trực tiếp mô hình embedding hoặc PhoBERT trên dữ liệu này; gán nhãn theo mệnh đề để xử
lý câu trộn lẫn; thêm nhãn riêng cho lời góp ý; thu thập phản hồi từ khoa hoặc trường khác để kiểm tra khả năng tổng
quát hóa.

\begin{thebibliography}{9}
\bibitem{text-vsfc} K. V. Nguyen, V. D. Nguyen, P. X. V. Nguyen, T. T. H. Truong, N. L.-T. Nguyen.
  \textit{UIT-VSFC: Vietnamese Students' Feedback Corpus for Sentiment Analysis}. 10th International Conference on
  Knowledge and Systems Engineering (KSE), 2018, tr. 19-24.
\bibitem{text-emb} AITeamVN. \textit{Vietnamese\_Embedding}. Hugging Face, \url{https://huggingface.co/AITeamVN/Vietnamese_Embedding}.
\end{thebibliography}
